
\section{Conclusion}
\label{sec:conclusion}

We proposed a tick-first reformulation of Standard Model structure within the HPA scan--projection paradigm.
In this framing, the only primitives are time as tick (scan iteration count) and CAP as the bounded-complexity closure rule, while finite observability appears through window projection and stability filtering (Section~\ref{sec:tick_calculus} and Section~\ref{sec:hpa_readout}).
At the CAP-selected anchor on the chosen $2$D screen, $(m,n)=(6,3)$, the provable folding core at $m=6$ compresses $64$ microstates to $21$ stable types with a canonical $18\oplus 3$ cyclic/boundary split (Section~\ref{sec:folding_core}).
Locality language is then introduced as a derived display structure by an addressing basis, whose minimality at the anchor is made explicit and auditable (Table~\ref{tab:addressing_selection}).
On the physical identification layer, we recorded interface hypotheses: gauge fields as defect-compensating connections; chirality as protocol selection among orientation classes; antimatter as conjugate readout under scan reversal; and mass/scale as time dictionaries in a Fibonacci log-time coordinate.
We provided auditable scripts and generated tables that reproduce the finite folding statistics, the addressing and chirality diagnostics, and the quantitative closures, and we stated falsifiable predictions formulated directly in the protocol language (Section~\ref{sec:falsifiability}).
In addition, we closed two main interface components at $(m,n)=(6,3)$: a unique field-level labeling of the $21$ stable types, and a closed mass-spectrum depth formula with a bounded-complexity rigidity signal.

\paragraph{Summary.}
\begin{itemize}
\item \textbf{Theorem-level anchor.} At $(m,n)=(6,3)$ the folding core yields $64\to 21$ and the canonical split $21=18\oplus 3$ (Section~\ref{sec:folding_core}).
\item \textbf{Closed interfaces.} We close a unique SM labeling map $\mathcal{L}_{\mathrm{SM}}$ (Theorem~\ref{thm:labeling_unique}) and a closed mass-depth template (Definition~\ref{def:mass_template}, Table~\ref{tab:mass_spectrum}).
\item \textbf{Audited normalizations and mixing.} Coupling/CP targets and mixing closures are recorded as bounded-complexity CAP selections with explicit audits (Section~\ref{sec:couplings_cp} and Appendix~\ref{app:closure_audit_details}).
\item \textbf{Falsifiability.} Protocol-level predictions and test channels are summarized in Section~\ref{sec:falsifiability}.
\item \textbf{Ledger.} A compact dependency and input ledger is recorded in Appendix~\ref{app:inference_ledger}.
\end{itemize}

\paragraph{Stability and AEC duality (interface).}
\InterfaceTag The stable-sector viewpoint used in this paper admits a control-theoretic dual dictionary: long-lived low-entropy structure can be modeled as predictive active error correction (AEC) that suppresses protocol-level mismatch relative to passive baselines while paying overhead and dissipation costs (Section~\ref{sec:introduction}; Appendix~\ref{app:isomorphism_dictionary}).

\paragraph{Closing (interface shorthand).}
\InterfaceTag In the tick-first framing, time is the update count, space is a derived addressing/locality dictionary, and matter/scale are protocol-stability and overhead dictionaries anchored at $m=6$.
In this sense, the matching-layer gaps recorded throughout (log mismatches for normalizations and $\Delta r$ depth shifts for masses) are the controlled slack between a closed template and an operational convention, and stable structure can be read as resistance against mismatch under finite readout.

